%% =====================================================================
%% 武汉大学本科毕业论文（设计）开题报告 —— LaTeX 模板
%% 依据原件「3-武汉大学本科毕业论文（设计）开题报告.doc」整理
%%
%% 编译： xelatex proposal.tex   （建议编译两次；需 XeLaTeX + ctex）
%% 说明： 只需修改下方「填写区」和文末六个 \proposalpart 的内容。
%%
%% 原件说明：第 5 行原为跨 4 列的大单元格，内含红色填写说明——
%%   「开题报告正文（应包含以下几方面内容：一、选题目的和意义；
%%     二、国内外研究现状和发展趋势；三、研究内容、研究方法、技术路线
%%     及可行性分析；四、项目特色与创新点；五、进度安排；六、主要参考文献。）」
%%   LaTeX 版将其拆成六个可跨页的小节，标题沿用原提示文字；
%%   红色说明本身不打印（保留在上面的注释里）。
%% =====================================================================
\documentclass[UTF8, a4paper, zihao=-4, fontset=fandol]{ctexart}
\usepackage{whu-form}

%% 原件页面：A4，上下 2.54cm、左右 3.17cm；框线约 0.5pt
\geometry{a4paper, top=2.54cm, bottom=2.54cm, left=3.17cm, right=3.17cm}
\setlength{\arrayrulewidth}{0.5pt}

%% ========================== 填写区 ==========================
\Title{基于深度学习的图像语义分割方法研究}
\StudentName{张三}
\StudentID{2020300000000}
\College{计算机学院}
\Major{计算机科学与技术}
\Advisor{李四}
\AdvisorTitle{教授}
% \showhints   % 取消注释：空白字段显示浅灰填写提示（定稿前请注释回去）
%% ============================================================

%% 列宽按原件 16.5 : 36.3 : 13.2 : 34.0 分配；行高 28.35pt(1cm)
\newlength{\prw}
\setlength{\prw}{\dimexpr\textwidth-8\tabcolsep-5\arrayrulewidth\relax}
\newlength{\pra}\setlength{\pra}{0.165\prw}
\newlength{\prb}\setlength{\prb}{0.363\prw}
\newlength{\prc}\setlength{\prc}{0.132\prw}
\newlength{\prd}\setlength{\prd}{0.340\prw}
\newlength{\prm}
\setlength{\prm}{\dimexpr\prb+\prc+\prd+4\tabcolsep+2\arrayrulewidth\relax}

\begin{document}
\whutitle{武汉大学本科毕业论文(设计)开题报告}

\begin{center}
\begin{tabular}{|p{\pra}|p{\prb}|p{\prc}|p{\prd}|}
\hline
%% 1 题目
\whulabel{28.35pt}{题\ 目}
& \multicolumn{3}{p{\prm}|}{\whuvaluecell{28.35pt}{\whufield{title}{（填写题目）}}} \\
\hline
%% 2 学生姓名 / 学号
\whulabel{28.35pt}{学生姓名}
& \whuvaluecell{28.35pt}{\whufield{studentname}{}}
& \whulabel{28.35pt}{学号}
& \whuvaluecell{28.35pt}{\whufield{studentid}{}} \\
\hline
%% 3 所在学院 / 专业
\whulabel{28.35pt}{所在学院}
& \whuvaluecell{28.35pt}{\whufield{college}{}}
& \whulabel{28.35pt}{专业}
& \whuvaluecell{28.35pt}{\whufield{major}{}} \\
\hline
%% 4 指导教师 / 职称
\whulabel{28.35pt}{指导教师}
& \whuvaluecell{28.35pt}{\whufield{advisor}{}}
& \whulabel{28.35pt}{职称}
& \whuvaluecell{28.35pt}{\whufield{advisortitle}{}} \\
\hline
\end{tabular}
\end{center}

\vspace{1.2em}
%% ---------- 开题报告正文：按原件六方面提示分节填写 ----------
%% 每个 \proposalpart 的第二个花括号内写内容；留空则自动留出填写空间。
\proposalpart{一、选题目的和意义}{}
\proposalpart{二、国内外研究现状和发展趋势}{}
\proposalpart{三、研究内容、研究方法、技术路线及可行性分析}{}
\proposalpart{四、项目特色与创新点}{}
\proposalpart{五、进度安排}{}
\proposalpart{六、主要参考文献}{}

\end{document}
